% Draft_13_unsteady_dg_proposed.tex
%
% SIDE NOTES for Draft_13.tex, written 2026-11-06. Not part of the paper.
% Standalone: compile with
%     pdflatex Draft_13_unsteady_dg_proposed.tex   (twice; Draft_13.aux must exist)
% References to the paper's own equations, figures, and sections resolve
% against Draft_13.aux through the xr package and print the paper's numbers.
% Everything numbered inside these notes carries an N prefix.
%
% Every number here is new simulation output awaiting review. None is in
% Draft_13.tex. Part III holds proposed paper text; nothing there is applied.
%
% Sources, all under trunk/Python_Simulations/Vector_Fields/VF_Robot/:
%   experiments/unsteady_dg/mc_unsteady_double_gyre.py  --tag n1000, --tag eps025_vmax008
%   experiments/unsteady_dg/mc_bent_double_gyre.py      --tag n500, --diag
%   experiments/unsteady_dg/mc_bickley_jet.py           --tag n1000, --geom
%   experiments/unsteady_dg/trench_offset.py, scope_ratio.py, path_ratio.py --tag n50
%   experiments/unsteady_dg/analyze.py, fig_notes.py (all four figures)
% Outputs: experiments/unsteady_dg/outputs/

\documentclass[11pt]{article}
\usepackage[margin=1in]{geometry}
\usepackage{amsmath,amssymb}
\usepackage{graphicx}
\usepackage{booktabs}
\usepackage{float}
\usepackage{algorithm}
\usepackage[font=small,labelfont=bf]{caption}
\usepackage{url}
\usepackage{xr}
\externaldocument{Draft_13}
\graphicspath{{figures/notes_unsteady/}}

\renewcommand{\theequation}{N\arabic{equation}}
\renewcommand{\thefigure}{N\arabic{figure}}
\renewcommand{\thetable}{N\arabic{table}}
\renewcommand{\thealgorithm}{N\arabic{algorithm}}
\setlength{\parskip}{4pt}

\newcommand{\Dtr}{$D$ tracker}
\newcommand{\Str}{$s_1$ tracker}

\title{What the $D$ and $s_1$ Trackers Track\\[2pt]
\large Notes on Draft 13 beyond the steady double gyre}
\author{Prepared for C. Waight}
\date{November 6, 2026}

\begin{document}
\maketitle

\noindent These notes collect one afternoon of new simulations on three
fields that remove the special structure of the steady double gyre. Part~I
summarizes what was found and what it means for the paper. Part~II gives
the evidence field by field. Part~III holds proposed paper text in the
paper's register, including old and new wording for four existing
sentences of Draft~13. The appendices record where every number came
from and which edits to Draft~13 were applied earlier in the session.
Equation, figure, and section numbers without an N prefix refer to
Draft~13.

\tableofcontents

%======================================================================
\part{Summary}
%======================================================================

\section{Why these runs}

Every simulated review of Draft~13 raised the same objection. The steady
double gyre is a benchmark on which both surrogate fields mark the
separatrix exactly, and a reviewer can fairly ask whether that is a
property of the trackers or of the benchmark. Section~\ref{sec:dg_example}
already lists the reasons. On the steady gyre the shear strain vanishes
everywhere, the vorticity vanishes on the separatrix, so
$\mathbf{J} = \mathbf{S}$ there and $D = -s_1^2$, the Hessian
$\mathbf{H}_D$ is diagonal, and the separatrix is a straight, stationary
line.

Three test fields remove these properties one group at a time. The
time-periodic double gyre adds shear and vorticity on the separatrix and
sets it moving, but keeps it straight. A bent double gyre, constructed for
these notes, keeps the flow steady but curves the separatrix and puts
strong vorticity on it. A Bickley jet replaces the gyres with a sheared
zonal jet whose cat's-eye eddies are bounded by curved separatrices on
which rotation exceeds strain. Each field runs a seven-start Monte
Carlo protocol modeled on Section~\ref{sec:dg_bench}, noise-free, with the
paper's formation, gains, and primitives unchanged.

\section{The short answer}

The two trackers do what they were designed to do. Each one rides a trench
of its own derived field: the \Dtr{} rides trenches of the Jacobian
determinant, and the \Str{} rides trenches of the smaller strain
eigenvalue. Neither field knows where the separatrix is. The trenches lie
on the separatrix only when the separatrix carries little rotation
compared with its strain, and the steady double gyre is the extreme case,
with no rotation on the separatrix at all.

As rotation on the separatrix grows, the trenches drift off it, and the
trackers drift with them (Fig.~\ref{fig:structures}). Straddle retention
against the true separatrix then falls in step with a single quantity the
cluster can compute from its own fit, the rotation-to-strain ratio
$|\hat\omega|/2\hat r$ at the centroid (Fig.~\ref{fig:scope}). Where its
median along the ride is below about 0.1, the \Dtr{} keeps the straddle in
97 to 100\% of trials. Near 0.3 both trackers keep it in a quarter to a
half of trials. Above 0.5 neither keeps it in more than 10\% of trials. The Santa Barbara record of
Section~\ref{sec:disc_ocean} sits at 0.15 to 0.20, between the regimes
that work and the regimes that partly work.

For the paper this is a framing problem, not a broken method. The
estimator, both primitives as trench followers, the noise comparison of
Section~\ref{sec:disc_noise}, and the objectivity result of
Section~\ref{sec:rotating} all stand. The sentences that tie the trackers
to separatrices in general need a scope condition, and the scope condition
can be stated with a number the robots measure themselves.

\section{What the trackers actually track}
\label{sec:what}

A separatrix is a property of the flow's topology. It is the streamline
that leaves one saddle and arrives at another, and in a steady flow it is
an exact transport boundary, since fluid never crosses a streamline. The
two surrogate fields are properties of the local velocity gradient
instead. Writing the fitted Jacobian through its mean, strain, and
rotation parts,
\begin{equation}
    \mathbf{J} =
    \begin{bmatrix}
        \mu + s_n & s_s - \omega/2 \\
        s_s + \omega/2 & \mu - s_n
    \end{bmatrix},
    \qquad
    D = \det\mathbf{J} = \mu^2 - r^2 + \tfrac{1}{4}\omega^2,
    \qquad
    s_1 = \mu - r,
    \label{eq:Jparts}
\end{equation}
with $r = \sqrt{s_n^2 + s_s^2}$ the strain rate. A trench of $D$ is a
curve where strain most outweighs rotation, the Okubo-Weiss partition read
as a landscape \cite{8,22}. A trench of $s_1$ is a curve of strongest
instantaneous compression or stretching, which includes the attracting and
repelling OECS of \cite{19}. Neither condition mentions streamlines.

On the steady double gyre the two descriptions agree because the vorticity
vanishes on the separatrix. There $\omega = 0$ and $\mu = 0$, so
$D = -r^2 = -s_1^2$, and the separatrix sits at the bottom of both
landscapes. Once the separatrix carries vorticity, the term
$\tfrac{1}{4}\omega^2$ in (\ref{eq:Jparts}) acquires a transverse gradient
and pushes the trench of $D$ sideways. This is the mechanism of
Section~\ref{sec:rotating}, where a rotating observer adds a uniform spin
$\Omega$ and the term $\Omega\omega$ of (\ref{eq:D_not_objective}) moves
the trench of $D$ off the separatrix. Here the flow supplies the spin
itself. The field $s_1$ contains no vorticity at all, and it marks curves
of the strain geometry. The direction of a streamline, however, depends on
the whole of $\mathbf{J}$, so where rotation enters $\mathbf{J}$ nothing
ties the streamline to the trenches of $s_1$.

Fig.~\ref{fig:structures} shows the result on all four fields. On the
steady gyre the three curves coincide. On the time-periodic gyre at peak
forcing both trenches sit within 0.035 of the separatrix. On the bent
gyre with $\delta = 0.2$ the trench of $D$ bows up to 0.078 past the
separatrix and the trench of $s_1$ up to 0.155, about twice the formation
radius $\rho = 0.075$. On the Bickley jet the trench of $s_1$ runs nearly
flat along the jet flank and crosses the U-shaped separatrix, and $D$ has
no trench near the separatrix at 40\% of positions along it.

\begin{figure}[t]
    \centering
    \includegraphics[width=\textwidth]{notes_structures.png}
    \caption{The separatrix (black) against the true trench of $D$
    (orange, solid) and of $s_1$ (blue, dashed) on the four test fields,
    computed from the analytic field rather than the fit. Light streamlines
    show the flow, the shaded regions are rotation-dominated ($D > 0$),
    crosses mark the saddles, and the dotted circle has the formation
    radius $\rho = 0.075$ for scale. (a) Steady double gyre, where all
    three curves coincide. (b) Time-periodic gyre, $\epsilon = 0.25$, at
    peak forcing $\sin\omega t = 1$. (c) Bent gyre, $\delta = 0.2$.
    (d) Bickley jet, with the eye center as a gray dot. Each trench is the
    nearest local minimum across the separatrix, so gaps mark heights or
    positions where none exists within the search window. Notice how far
    the trenches bow in (c) and how the $s_1$ trench in (d) crosses the
    separatrix instead of following it.}
    \label{fig:structures}
\end{figure}

What the trackers follow is therefore meaningful, but it is a different
object from the one the paper names. The case for the trackers is that
$s_1$ trenches mark where floating material is attracted or repelled most
strongly right now, and floating material has been observed at sea to
collect onto attracting OECS \cite{12}. Furthermore, in unsteady flow
streamlines are not particle paths, so the instantaneous separatrix is not
a material boundary anyway. The case against
is that on a steady flow the separatrix is the transport boundary, and a
trench that leaves it is a curve that fluid crosses. Both cases are
correct. The paper can keep its claims by naming the trench as the target
and stating when the trench marks the separatrix.

\section{How bad, in numbers}

Table~\ref{tab:pooled} gives the pooled outcomes on every field, and
Fig.~\ref{fig:scope} plots them against the median fitted
$|\hat\omega|/2\hat r$ along each ride. Three regimes appear.

\begin{itemize}
\item Below about 0.1 (steady gyre, time-periodic gyre at
$\epsilon = 0.1$), both trackers reach the downstream saddle in more than
96\% of trials. The \Dtr{} keeps the straddle in 97 to 100\%, and the
\Str{} in 59 to 100\%.
\item Near 0.3 (bent gyre, $\delta = 0.1$), success holds or falls modestly
(71\% for the \Dtr{}, 99.7\% for the \Str{}), while retention falls to
53\% and 24\%.
\item Above 0.5 (bent gyre at $\delta = 0.2$, Bickley jet), retention is
under 10\% for both, and on the Bickley jet neither tracker reaches the
downstream saddle in more than 4.1\% of trials.
\end{itemize}

The time-periodic gyre at $\epsilon = 0.25$ fails for a different reason.
Its separatrix moves faster than either tracker's across-trench channel
can command, and doubling $c_{\max}$ recovers most of the loss
(Section~\ref{sec:tp}). It is drawn hollow in Fig.~\ref{fig:scope} and
left out of the regimes above.

Success is a weaker signal than retention. A trench that leaves the
separatrix in mid-span still meets it near the saddles on the gyres, so a
tracker can lose the separatrix and still arrive. Retention measures
whether the cluster stayed on the separatrix, which is what the paper
claims.

\section{What survives and what changes}

The estimator, the minimality result, the sensitivity analysis, both
surrogate fields, both primitives, the noise comparison, and the
objectivity comparison are untouched. They are statements about the fit
and about trench following, and none of the new runs contradicts them.

Four passages tie the trackers to separatrices in general and need a scope
condition. Part~III gives old and new wording for each.
\begin{enumerate}
\item Section~\ref{sec:surrogates}: ``In the flows studied here,
separatrices run along the trenches of $D$ inside the strain-dominated
regions.''
\item The abstract: ``Each tracker locates the nearest separating
boundary and navigates to and along it.''
\item Section~\ref{sec:disc_limits}: the $s_1$ tracker's ``passage through
the origin rests on the benchmark's vanishing shear strain,'' which the
time-periodic gyre now tests directly.
\item Section~\ref{sec:disc_limits} gains one paragraph on the scope
condition and on the ratio as an online check.
\end{enumerate}

One addition could become a contribution. The cluster computes
$|\hat\omega|/2\hat r$ from the first-order coefficients of the same fit,
the least noisy ones (Section~\ref{sec:sensitivity}), so it could flag
online when its trench is unlikely to mark a transport boundary.
Algorithm~\ref{alg:indicator} states the computation. It has not been
built into either primitive.

\section{Open items}

\begin{enumerate}
\item \textbf{Score each tracker against its own true trench.} The
retention numbers here are measured against the separatrix. The same
trials scored against the true trench of each tracker's own field would
separate ``wrong target'' from ``lost target.'' Fig.~\ref{fig:rides}
suggests wrong target for the \Str{} on the Bickley jet, since its rides
end on its own flat trench. The \Dtr{} failures from S6 and S7 on the bent
gyre are not yet explained.
\item \textbf{Threshold or trend.} Five steady-structure conditions in
three families give a monotone trend, not a threshold. Intermediate
values of $\delta$ and $\epsilon$, and one or two more field families,
would show whether a usable cutoff exists.
\item \textbf{Noise.} Every run here is noise-free. The fitted ratio uses
first-order coefficients, so it should degrade slowly with noise, but that
is untested.
\item \textbf{Lagrangian reference.} On the time-periodic gyre the
material boundary is the FTLE ridge of the unsteady flow, not the
instantaneous separatrix. No comparison against it was run.
\item \textbf{Bickley citation.} The paper-search corpus holds no primary
source for the Bickley jet, so Part~III cites none. A source would need
to be checked and added before the subsection could go in.
\item \textbf{Constructed field.} The bent gyre was built for these notes
and has no literature source. If it goes in the paper it should be
described as constructed.
\end{enumerate}

%======================================================================
\part{Evidence}
%======================================================================

\section{Common protocol}
\label{sec:protocol}

Every field uses the paper's pentagon-plus-center formation with
$\rho = 0.075$, the double-gyre operating point of Table~\ref{tab:params},
and the two primitives exactly as in Sections~\ref{sec:sep_controller}
and \ref{sec:s1_controller}. The runs are noise-free. Each trial draws a
random heading, and on the time-periodic gyre a random forcing phase, and
the two trackers share both draws, so every comparison between them is
paired. Algorithm~\ref{alg:trial} gives one trial.

The scoring follows Section~\ref{sec:dg_bench} with two changes. First,
straddle, acquisition, and success are measured against the moving or
curved separatrix and saddles instead of $x = 0$ and fixed points. Second,
one stopping rule serves both trackers, which answers the reviewer comment
on inconsistent criteria. After contact with the downstream saddle the run
continues for 100 steps, one forcing period of the time-periodic gyre, and
the hold is the fraction of those steps spent within 0.06 of it.

\begin{algorithm}[H]
\caption{One Monte Carlo trial on a test field}
\label{alg:trial}
\begin{tabbing}
\hspace{1.5em}\=\hspace{1.5em}\=\hspace{1.5em}\=\kill
\textbf{Input:} field $\mathbf{v}(\mathbf{p}, t)$; separatrix $\Gamma(t)$;
downstream saddle $\mathbf{p}^*_1(t)$; upstream saddle $\mathbf{p}^*_2(t)$;
start $\mathbf{p}_0$; primitive $\pi \in \{D, s_1\}$\\
Draw heading $\theta_0 \sim U[0, 2\pi)$; on the time-periodic gyre draw
$\omega t_0 \sim U[0, 2\pi)$, else $t_0 \leftarrow 0$\\
Place the formation at $\mathbf{p}_0$ with heading $\theta_0$;
$t \leftarrow t_0$; outcome $\leftarrow$ timeout\\
\textbf{Loop} for $k = 1, \ldots, 500$:\\
\> Sample $\mathbf{v}$ at the six robots at time $t$, fit, and command
$\dot{\mathbf{p}}_c \leftarrow \pi(\text{fit})$\\
\> Move the robots one control period; $t \leftarrow t + \Delta t$\\
\> Record whether robots lie on both sides of $\Gamma(t)$ (straddle)\\
\> \textbf{If} $\|\mathbf{p}_c - \mathbf{p}^*_1(t)\| < 0.06$:
outcome $\leftarrow$ success; \textbf{break}\\
\> \textbf{If} $\|\mathbf{p}_c - \mathbf{p}^*_2(t)\| < 0.06$:
outcome $\leftarrow$ reversal; \textbf{break}\\
\> \textbf{If} $\mathbf{p}_c$ leaves the domain box:
outcome $\leftarrow$ exit; \textbf{break}\\
\textbf{If} outcome is success: run 100 more steps and record the
fraction within 0.06 of $\mathbf{p}^*_1(t)$ (hold)\\
Retained $\leftarrow$ success \textbf{and} straddle held at every step
from the first straddle to contact
\end{tabbing}
\end{algorithm}

Three controls check the harness. At $\epsilon = 0$ the time-periodic
runner and at $\delta = 0$ the bent runner both reproduce the paper's
noise-free result, with 100\% success, retention, and hold for both
trackers from every start. The bent field at $\delta = 0$ equals the
paper's field to machine precision. The first 200 trials of each 1000-trial
time-periodic cell repeat an earlier 200-trial run exactly in every
column.

\begin{table}[t]
\centering
\caption{The test fields. The separatrix vorticity is its largest
magnitude along the separatrix. Trench offsets are the largest distance
from the separatrix to the nearest local minimum of the true $D$ or $s_1$
across it, horizontal on the gyres (excluding $|y| < 0.1$) and vertical on
the jet; for the time-periodic gyre they are taken at peak forcing.}
\label{tab:fields}
\small
\begin{tabular}{p{0.2\textwidth}p{0.22\textwidth}p{0.13\textwidth}p{0.17\textwidth}p{0.11\textwidth}}
\toprule
Field & Removes & Separatrix vorticity & $D$ trench offset & $s_1$ trench offset \\
\midrule
Steady gyre                 & nothing                       & 0          & 0     & 0     \\
Time-periodic, $\epsilon=0.1$  & zero shear, stationarity   & 0.063      & 0.019 & 0.019 \\
Time-periodic, $\epsilon=0.25$ & zero shear, stationarity   & 0.157      & 0.035 & 0.035 \\
Bent, $\delta=0.1$           & straightness, $\omega = 0$ on it & 0.60    & 0.041 & 0.085 \\
Bent, $\delta=0.2$           & straightness, $\omega = 0$ on it & 1.19    & 0.078 & 0.155 \\
Bickley jet                  & all of the above             & 0.82 to 1.96 & absent at 40\% of positions, else $\le 0.148$ & 0.133 \\
\bottomrule
\end{tabular}
\end{table}

\begin{table}[t]
\centering
\caption{Pooled outcomes, noise-free. Rates in percent over all trials
of the condition; hold is the median over successful trials. Trials:
$7 \times 10^3$ per tracker on the time-periodic gyre and the Bickley jet,
$3.5 \times 10^3$ on the bent gyre. The row marked $2c_{\max}$ is a
diagnostic with $c_{\max} = 0.08$.}
\label{tab:pooled}
\begin{tabular}{llrrrrrr}
\toprule
Field & Tracker & Success & Retention & Reversal & Exit & Timeout & Hold \\
\midrule
Steady gyre & $D$ & 100 & 100 & 0 & 0 & 0 & 1.00 \\
            & $s_1$ & 100 & 100 & 0 & 0 & 0 & 1.00 \\
Time-periodic, $\epsilon=0.1$ & $D$ & 97.3 & 97.3 & 2.7 & 0 & 0 & 1.00 \\
            & $s_1$ & 96.2 & 59.1 & 3.8 & 0 & 0 & 1.00 \\
Time-periodic, $\epsilon=0.25$ & $D$ & 30.0 & 8.6 & 16.1 & 53.9 & 0 & 0.39 \\
            & $s_1$ & 59.4 & 15.5 & 28.0 & 12.6 & 0 & 0.27 \\
\quad same, $2c_{\max}$ & $D$ & 62.4 & 46.3 & 18.4 & 19.2 & 0 & 0.89 \\
            & $s_1$ & 61.7 & 28.0 & 25.5 & 12.8 & 0 & 1.00 \\
Bent, $\delta=0.1$ & $D$ & 71.4 & 52.5 & 0 & 28.6 & 0 & 1.00 \\
            & $s_1$ & 99.7 & 23.5 & 0.3 & 0 & 0 & 1.00 \\
Bent, $\delta=0.2$ & $D$ & 62.1 & 9.7 & 0 & 37.9 & 0 & 1.00 \\
            & $s_1$ & 53.4 & 0.5 & 9.6 & 28.8 & 8.2 & 1.00 \\
Bickley jet & $D$ & 4.1 & 0 & 9.3 & 55.9 & 30.7 & 1.00 \\
            & $s_1$ & 0 & 0 & 0 & 0.3 & 99.7 & n/a \\
\bottomrule
\end{tabular}
\end{table}

\section{Time-periodic double gyre}
\label{sec:tp}

\subsection{Field}

The time-periodic gyre is the unsteady member of the family of \cite{17},
the field the paper already uses with $\epsilon = 0$. It replaces $x_f$ in
(\ref{eq:dg_u}) by
$f(x_f, t) = \epsilon\sin(\omega t)\,x_f^2 + (1 - 2\epsilon\sin\omega t)\,x_f$
and multiplies $v$ by $\partial f/\partial x_f$, with $\omega = 2\pi/10$.
The instantaneous separatrix, the line $u = 0$, is where $f = 1$. It
stays straight but swings sideways, with amplitude 0.099 at
$\epsilon = 0.1$ and 0.236 at $\epsilon = 0.25$ and peak speeds of 0.063
and 0.157. The saddles ride along the walls with it, and the flow on it
still runs downward, so $\mathbf{p}^*_1(t)$ stays the downstream saddle.

On the separatrix $\sin\pi f = 0$, which removes every term of the
gradient except the one carrying $\partial^2 f/\partial x_f^2 =
2\epsilon\sin\omega t$. The vorticity there is therefore
$-2\pi A\epsilon\sin(\omega t)\sin(\pi y_f)$ and the shear strain half of
it, with peaks $2\pi A\epsilon$ and $\pi A\epsilon$. Both vanish on the
steady gyre. At the crest, where $u_x$ also vanishes, the ratio
$|\omega|/2r$ is exactly 1 whenever $\sin\omega t \neq 0$, while the
median along the separatrix is only 0.018 at $\epsilon = 0.1$
(Table~\ref{tab:ratio}). Rotation matters at a single point near the
crest and nowhere else.

\subsection{Results at $\epsilon = 0.1$}

Both trackers reach $\mathbf{p}^*_1(t)$ in more than 96\% of trials and
hold it for the full period in every successful trial, since the saddle
moves at most 0.063 along the wall, below the $k\,c_{\max} = 0.12$ each
component of either hold law can command. The \Dtr{} keeps the straddle
in every trial that succeeds. Its median tracking error of 0.017 is close
to the 0.015 to 0.019 offset of its own trench at peak forcing, so the
\Dtr{} follows its trench closely and the trench sits close to the
separatrix.

The \Str{} reaches the saddle as often but keeps the straddle in only
59\% of trials, and in 9 to 14\% from the three starts above the origin
(Table~\ref{tab:perstart}). Fig.~\ref{fig:crest} shows where it loses it.
Of the 2595 successful trials that lose the straddle, 97\% lose it within
0.1 of $y = 0$. Near the crest the strain on the separatrix falls to at
most $\pi A\epsilon = 0.031$, against $\pi^2 A = 0.99$ at the saddles, so
$s_1$ is nearly flat across the trench. The across-trench command of
(\ref{eq:s1_law}) scales with the slope of $s_1$, and a flat trench gives
it too little to keep pace with a moving one. The \Dtr{}'s across-trench
command is a Newton step, a distance, which the flattening does not
weaken. The steady gyre hides this weakness because its separatrix does
not move.

\begin{figure}[t]
    \centering
    \includegraphics[width=0.62\textwidth]{notes_crest_loss.png}
    \caption{Time-periodic gyre, $\epsilon = 0.1$. Height of the centroid
    at the first step on which the \Str{} loses the straddle, over the 2595
    successful trials that lose it. The shaded band is $|y| < 0.1$ around
    the crest, which holds 97\% of the losses. The \Dtr{} loses the
    straddle in no successful trial.}
    \label{fig:crest}
\end{figure}

Reversals at $\epsilon = 0.1$ come only from S1 and S2, and for the
\Str{} also from S4, in 7 to 11\% of their trials. None of them rides down
first, and each reaches $\mathbf{p}^*_2(t)$ in a median of 17 to 20 steps.
Every reversing trial starts at a phase where the displaced separatrix is
nearly as far from the start as the top-wall trench, with a median ratio
of the two distances of 0.98 and a minimum of 0.89, against 0.47 in trials
that do not reverse. The cluster acquires the wrong trench rather than
turning on the right one. This is the behavior the abstract claims, since
the tracker goes to the nearest trench, but on the time-periodic gyre the
nearest trench is sometimes the wall's.

\subsection{Results at $\epsilon = 0.25$}

At $\epsilon = 0.25$ the separatrix moves at up to 0.157, faster than the
0.12 cap on either tracker's across-trench channel. The \Dtr{} succeeds in
30\% of trials and the \Str{} in 59\%, and neither holds the saddle, which
outruns both hold laws along the wall. Reversals rise to 16\% and 28\%, and
in 91\% of the \Dtr{}'s reversals and 86\% of the \Str{}'s the top-wall
trench was nearer to the start than the displaced separatrix. The
\Dtr{}'s remaining failures, 54\% of its trials, leave the basin at a
bottom corner after missing the moving saddle.

A diagnostic run with $c_{\max}$ doubled to 0.08, which lifts the cap to
0.24, separates the speed limit from the structure. The \Dtr{}'s success
rises to 62\% and its hold to 0.89 of the period, and the \Str{}'s success
to 62\% with a full hold. Reversals stay at 18\% and 26\%, since they come
from the start geometry rather than the speed limit. The $\epsilon = 0.25$
failures therefore belong to the operating point, not to the surrogate
fields, and Fig.~\ref{fig:scope} draws them hollow for that reason.

\begin{table}[t]
\centering
\caption{Success and straddle retention by start, in percent. Starts
S1 to S7 are those of Fig.~\ref{fig:traverse_vs_logic_c}. The Bickley
starts B1 to B7 are on the separatrix at $x = -0.3$, 0, and 0.3, then off
it by $2\rho$ above and below at $x = -0.25$ and by $1.3\rho$ above and
below at $x = 0.15$.}
\label{tab:perstart}
\small
\begin{tabular}{lll rrrrrrr}
\toprule
Field & Tracker & Metric & S1 & S2 & S3 & S4 & S5 & S6 & S7 \\
\midrule
TP, $\epsilon=0.1$ & $D$ & success & 90.7 & 90.7 & 100 & 100 & 100 & 100 & 100 \\
 & & retention & 90.7 & 90.7 & 100 & 100 & 100 & 100 & 100 \\
 & $s_1$ & success & 88.9 & 93.0 & 100 & 91.5 & 100 & 100 & 100 \\
 & & retention & 9.1 & 9.8 & 13.6 & 81.4 & 100 & 100 & 100 \\
\midrule
TP, $\epsilon=0.25$ & $D$ & success & 28.1 & 22.7 & 29.2 & 52.7 & 26.6 & 27.3 & 23.6 \\
 & & retention & 0 & 0 & 0 & 0 & 17.1 & 23.0 & 19.9 \\
 & $s_1$ & success & 50.0 & 61.1 & 27.0 & 49.3 & 59.0 & 88.5 & 80.7 \\
 & & retention & 0 & 0 & 0 & 0 & 26.6 & 41.4 & 40.7 \\
\midrule
Bent, $\delta=0.1$ & $D$ & success & 100 & 100 & 100 & 100 & 100 & 0 & 0 \\
 & & retention & 69.6 & 44.8 & 55.2 & 97.8 & 100 & 0 & 0 \\
 & $s_1$ & success & 100 & 100 & 100 & 98.2 & 100 & 100 & 100 \\
 & & retention & 0 & 0 & 0 & 0.8 & 49.6 & 49.4 & 65.0 \\
\midrule
Bent, $\delta=0.2$ & $D$ & success & 100 & 100 & 100 & 100 & 34.4 & 0 & 0 \\
 & & retention & 0 & 0 & 0 & 33.2 & 34.4 & 0 & 0 \\
 & $s_1$ & success & 96.4 & 55.6 & 94.6 & 91.8 & 8.4 & 1.4 & 25.8 \\
 & & retention & 0 & 0 & 0 & 0 & 0 & 0 & 3.2 \\
\midrule
 & & & B1 & B2 & B3 & B4 & B5 & B6 & B7 \\
Bickley jet & $D$ & success & 0 & 3.5 & 25.1 & 0 & 0 & 0 & 0 \\
 & $s_1$ & success & 0 & 0 & 0 & 0 & 0 & 0 & 0 \\
\bottomrule
\end{tabular}
\end{table}

\section{Bent double gyre}
\label{sec:bent_notes}

\subsection{Field}

The time-periodic gyre keeps its separatrix straight. The bent gyre curves
it while staying steady and incompressible, by bending the stream function
of the paper's gyre,
\begin{equation}
    \psi = A\sin(\pi g)\sin(\pi y_f), \qquad
    g = x_f + \delta\sin(\pi y_f), \qquad
    u = -\frac{\partial\psi}{\partial y}, \quad
    v = \frac{\partial\psi}{\partial x}.
    \label{eq:bent}
\end{equation}
At $\delta = 0$ it is (\ref{eq:dg_u}) exactly. For $\delta > 0$ the
separatrix is the streamline $g = 1$, the curve $x = -\delta\cos(\pi y)$,
bowed by $\delta$ at $y = 0$ between the unchanged saddles $(0, \pm 0.5)$.
Two checks confirm the construction. The flow on the curve is tangent to
it, with $v/u$ equal to its slope at every point tested, and it runs
downward, so $\mathbf{p}^*_1 = (0, -0.5)$ stays the downstream saddle.
The bend was chosen for these notes and has no literature source.

The vorticity on the separatrix reaches 0.60 at $\delta = 0.1$ and 1.19 at
$\delta = 0.2$, comparable to the strain $\pi^2 A = 0.99$ at the saddles.
Both trenches bow further than the separatrix toward its concave side and
meet it near the saddles (Fig.~\ref{fig:structures}(c)). The $s_1$ trench
lies up to 0.085 off at $\delta = 0.1$, already more than $\rho$, and up
to 0.155 off at $\delta = 0.2$. The $D$ trench lies up to 0.041 and 0.078
off.

\subsection{Results}

From the four starts at or above the origin, the \Dtr{} reaches
$\mathbf{p}^*_1$ in every trial at both bends, and the \Str{} in 98 to
100\% at $\delta = 0.1$ and 56 to 96\% at $\delta = 0.2$
(Table~\ref{tab:perstart}). Both hold the saddle in every successful trial
at $\delta = 0.1$. Retention falls for both, and furthest for the \Str{},
whose trench lies further from the separatrix. At $\delta = 0.1$ the
\Str{} keeps the straddle in none of the trials from S1 to S3, while still
reaching the saddle in all of them. Its rides in Fig.~\ref{fig:rides}(f)
swing well off the separatrix through mid-height and rejoin it near the
downstream saddle.

From S6 and S7 the \Dtr{} fails in every trial at both bends. In
Fig.~\ref{fig:rides}(b) and (c) these rides descend to the bottom wall and
ride the bottom-wall trench east until they leave the basin. The bend
moves the separatrix away from these two starts. At $\delta = 0.1$ they
sit 0.24 and 0.18 from it, $3.2\rho$ and $2.4\rho$, against $2\rho$ and
$1.3\rho$ on the steady gyre, and further still from the bowed $D$ trench.
Why the \Dtr{} takes the wall trench from there, while the \Str{} still
succeeds from all three lower starts at $\delta = 0.1$, is not yet
diagnosed.

\begin{figure}[p]
    \centering
    \includegraphics[width=\textwidth]{notes_rides.png}
    \caption{One ride per start with heading 0 (and forcing phase 0 on the
    time-periodic gyre), \Dtr{} on top and \Str{} below. Gray dots are
    starts. A filled end marker reached the downstream saddle; a hollow one
    did not. Black curves are separatrices, and on the time-periodic gyre
    the shaded band is its sweep with the dotted line its mean. Thin
    colored curves on the bent gyre and the jet are the tracker's own true
    trench. (a, e) Time-periodic, $\epsilon = 0.1$: all fourteen rides
    succeed, and the \Str{} rides swing across the band near the crest.
    (b, c, f, g) Bent gyre: the \Dtr{} rides from S6 and S7 descend to the
    bottom wall and leave the basin along it, and the \Str{} rides follow
    its bowed trench. (d, h) Bickley jet: the \Dtr{} rides scatter, and the
    \Str{} rides end on its nearly flat trench below the saddles. These
    single rides illustrate the mechanisms; the rates are in
    Tables~\ref{tab:pooled} and \ref{tab:perstart}.}
    \label{fig:rides}
\end{figure}

\section{Bickley jet}
\label{sec:bickley_notes}

\subsection{Field}

The gyres have separatrices on which strain dominates over most of their
length. A Bickley jet is a sheared zonal jet whose traveling waves roll the
flow on each flank into a row of cat's-eye eddies. Seen in the frame of
its wave, a single-mode jet is steady,
\begin{equation}
    \psi = -UL\tanh(y/L) + c\,y
        + UL\epsilon\,\mathrm{sech}^2(y/L)\cos k(x + 0.5),
    \label{eq:bickley}
\end{equation}
with $u = -\partial\psi/\partial y$ and $v = \partial\psi/\partial x$. The
parameters $U = 0.3$, $L = 0.2$, $c = 0.3U$, $\epsilon = 0.3$, and
$k = 2\pi$ match the gyre benchmark in speed (about 0.3) and separatrix
length (about 1). The northern eye centered at $(0, 0.172)$ is bounded
below by a separatrix that runs east from the saddle at $(-0.5, 0.292)$ to
the saddle at $(0.5, 0.292)$, dipping 0.26 below them in a U, with
maximum slope 0.73 and length 1.14. The eye is 0.28 tall at $x = 0$, so
the formation fits inside it.

The separatrix is traced numerically along the unstable manifold of the
upstream saddle. It ends 0.0009 from the downstream saddle, $\psi$ stays
constant along it to $1.5 \times 10^{-13}$, and the flow on it is tangent
and runs east. On it the vorticity ranges from 0.82 to 1.96, against a
strain rate of 0.33 at the saddles, so rotation exceeds strain over most
of its length (Table~\ref{tab:ratio}).

\subsection{Results}

Neither tracker follows this separatrix. Over $10^3$ trials from each of
the seven starts, the \Dtr{} reaches the downstream saddle in 4.1\% of
trials, all from B2 and B3, and the \Str{} in none. Neither keeps the
straddle in any trial. The \Dtr{} leaves the domain in 56\% of trials and
reaches the upstream saddle in 9\%. The \Str{} runs the full 500 steps
without reaching either saddle in 99.7\% of trials.

Fig.~\ref{fig:structures}(d) and Fig.~\ref{fig:rides}(h) together explain
the \Str{}. Its trench near the separatrix runs nearly flat along the jet
flank near $y = 0.15$, crossing the U-shaped separatrix twice and passing
0.13 to 0.14 below the saddles. In the full runs the median final height
is 0.16 from every start except B5, which starts below the jet core, and
in Fig.~\ref{fig:rides}(h) the rides from the other starts end on that
trench. The \Str{} follows its own trench, and the trench is not the
separatrix. The \Dtr{}'s rides scatter in Fig.~\ref{fig:rides}(d), as
expected where $D$ has no trench near the separatrix at 40\% of
positions along it, and their individual mechanisms are not diagnosed.

\section{The scope indicator}
\label{sec:scope_notes}

\subsection{Along the true separatrix}

Table~\ref{tab:ratio} reports the true ratio $|\omega|/2r$ along each
separatrix, weighted by arc length. A ratio below 1 means strain
dominates, which in incompressible flow is $D < 0$. The median orders the
fields in the same sequence as their retention, from 0 on the steady gyre
to 0.79 on the Bickley jet. The maximum is a poor summary on the gyres,
since it reaches exactly 1 at the crest of every perturbed gyre. Only the
Bickley jet has stretches of separatrix that are rotation-dominated, 40\%
of its length.

\begin{table}[t]
\centering
\caption{Rotation-to-strain ratio $|\omega|/2r$ along the true separatrix
(arc-length weighted; the time-periodic gyre averaged over 40 phases), and
the median over trials of the fitted $|\hat\omega|/2\hat r$ at the
centroid along each ride (50 trials per start), with retention from
Table~\ref{tab:pooled}. The Santa Barbara value is the fitted median along
the tracker paths from 7 to 21~h in Section~\ref{sec:disc_ocean}.}
\label{tab:ratio}
\small
\setlength{\tabcolsep}{4pt}
\begin{tabular}{lrrrr rr rr}
\toprule
 & \multicolumn{4}{c}{True, along separatrix} & \multicolumn{2}{c}{Fitted, along ride} & \multicolumn{2}{c}{Retention (\%)} \\
\cmidrule(lr){2-5}\cmidrule(lr){6-7}\cmidrule(lr){8-9}
Field & median & 90th pct. & $>0.5$ & $>1$ & $D$ & $s_1$ & $D$ & $s_1$ \\
\midrule
Steady gyre                & 0.000 & 0.000 & 0\%    & 0\%    & 0.007 & 0.027 & 100  & 100 \\
Time-periodic, $\epsilon=0.1$  & 0.018 & 0.126 & 2.2\%  & 0\% & 0.083 & 0.083 & 97.3 & 59.1 \\
Time-periodic, $\epsilon=0.25$ & 0.042 & 0.285 & 5.2\%  & 0\% & 0.055 & 0.236 & 8.6  & 15.5 \\
Bent, $\delta=0.1$         & 0.250 & 0.673 & 14.5\% & 0\%    & 0.257 & 0.327 & 52.5 & 23.5 \\
Bent, $\delta=0.2$         & 0.460 & 0.862 & 39.7\% & 0\%    & 0.517 & 0.578 & 9.7  & 0.5 \\
Bickley jet                & 0.786 & 2.976 & 100\%  & 39.5\% & 0.591 & 0.569 & 0    & 0 \\
\midrule
Santa Barbara record       & & & & & \multicolumn{2}{c}{0.15 to 0.20} & & \\
\bottomrule
\end{tabular}
\end{table}

\subsection{Along the cluster's own ride}

The true ratio along the separatrix is not available to a cluster that
does not know where the separatrix is. The fitted ratio at the centroid
is. It uses only the first-order coefficients of the fit,
Algorithm~\ref{alg:indicator}, and it is the quantity
Section~\ref{sec:disc_ocean} already reports for the ocean record. The fitted medians track the true ones on
the steady fields, 0.257 against 0.250 on the bent gyre at $\delta = 0.1$,
and they run higher on the time-periodic gyre, where the rides spend time
off the separatrix near the crest. The $\epsilon = 0.25$ \Dtr{} value is
low because most of its trials leave the separatrix early and ride
elsewhere.

\begin{algorithm}[H]
\caption{Online scope indicator from the cluster's own fit}
\label{alg:indicator}
\begin{tabbing}
\hspace{1.5em}\=\hspace{1.5em}\=\hspace{1.5em}\=\kill
\textbf{Input:} fitted first-order coefficients
$\hat a_2, \hat a_3, \hat b_2, \hat b_3$ at the centroid, each step\\
$\hat s_n \leftarrow (\hat a_2 - \hat b_3)/2$; \quad
$\hat s_s \leftarrow (\hat a_3 + \hat b_2)/2$; \quad
$\hat r \leftarrow \sqrt{\hat s_n^2 + \hat s_s^2}$\\
$\hat\omega \leftarrow \hat b_2 - \hat a_3$\\
$q_\omega \leftarrow |\hat\omega| / 2\hat r$ \quad (rotation-to-strain ratio)\\
Keep a running median of $q_\omega$ over the ride\\
\textbf{If} the running median exceeds a threshold $q^*$: report that the
trench being ridden is unlikely to mark a separatrix\\
\textbf{Note:} $q^*$ is not chosen. Fig.~\ref{fig:scope} suggests a value
between 0.1 and 0.3, and the open items above list what would fix it.
\end{tabbing}
\end{algorithm}

\begin{figure}[t]
    \centering
    \includegraphics[width=\textwidth]{notes_scope.png}
    \caption{Straddle retention (left) and downstream-saddle success
    (right) against the median fitted $|\hat\omega|/2\hat r$ along the
    ride, for the \Dtr{} (orange circles) and the \Str{} (blue triangles).
    Ratios come from 50 reruns per start, rates from the full runs of
    Table~\ref{tab:pooled}. Hollow markers are the time-periodic gyre at
    $\epsilon = 0.25$, which fails on the speed limit rather than on
    rotation. The shaded band is the Santa Barbara record. Notice that
    retention falls steadily with the ratio for both trackers, while
    success holds longer, because the trenches on the gyres rejoin the
    separatrix near the saddles.}
    \label{fig:scope}
\end{figure}

\subsection{What the ocean number means}

The Santa Barbara record's fitted median of 0.15 to 0.20 sits between the
time-periodic gyre at $\epsilon = 0.1$ (0.083) and the bent gyre at
$\delta = 0.1$ (0.26 to 0.33). On the synthetic fields that range keeps
the straddle in roughly half or more of the \Dtr{}'s trials. The ocean
comparison is not one to one. Its reference is an FTLE ridge, not a
separatrix, its formation spans five grid cells of interpolated radar
data, and its trackers ran at a different operating point. The
agreement with the FTLE ridge in Section~\ref{sec:disc_ocean} is
therefore consistent with these results, and it sits in the range where
the synthetic fields predict partial rather than guaranteed agreement.

%======================================================================
\part{Proposed paper text}
%======================================================================

\noindent The text below is in the paper's register and uses the paper's
labels, which resolve here against Draft~13. Equations and tables carry
N numbers in these notes and would take the paper's numbering on
insertion. With all three subsections added the paper compiles at 13
pages.

\section{New subsections for Section~\ref{sec:dg_bench}}

\subsection*{IV-E \quad Time-Periodic Double Gyre}

The steady benchmark has zero shear strain everywhere, a diagonal
$\mathbf{H}_D$, and a stationary separatrix on which
$\mathbf{J} = \mathbf{S}$. The time-periodic member of the family of
\cite{17} removes all three. It replaces $x_f$ in (\ref{eq:dg_u}) by
\begin{equation}
    f(x_f, t) = \epsilon \sin(\omega t)\, x_f^2
        + \bigl(1 - 2\epsilon \sin(\omega t)\bigr) x_f
    \label{eq:dg_f}
\end{equation}
and multiplies $v$ by $\partial f/\partial x_f$, with $\omega = 2\pi/10$.
The instantaneous separatrix, the line $u = 0$, stays straight at the $x$
where $f = 1$, but it oscillates sideways with amplitude 0.099 at
$\epsilon = 0.1$ and 0.236 at $\epsilon = 0.25$, at peak speeds of 0.063
and 0.157. The saddles move with it along the walls, and the flow on it
still runs downward, so $\mathbf{p}^*_1(t)$ remains the downstream saddle.
On the separatrix the vorticity reaches $2\pi A\epsilon$ and the shear
strain $\pi A\epsilon$, so $\mathbf{J} \neq \mathbf{S}$ there. At peak
forcing the true trenches of $D$ and $s_1$ sit up to 0.019
($\epsilon = 0.1$) and 0.035 ($\epsilon = 0.25$) off the separatrix, on
the side toward which it has moved.

Both trackers run from the seven starts of
Fig.~\ref{fig:traverse_vs_logic_c} at $\epsilon = 0$, 0.1, and 0.25,
noise-free, with $10^3$ trials per start, tracker, and $\epsilon$. Each
trial draws its heading and its initial phase $\omega t_0$ uniformly from
$[0, 2\pi)$, and the two trackers share both. Straddle, acquisition, and
success are scored against the moving separatrix $x_s(t)$ and saddles. A
trial ends at contact with either instantaneous saddle (centroid within
0.06), when the centroid leaves the box $|x| < 1$, $|y| < 0.6$, or after
500 steps. After contact with $\mathbf{p}^*_1(t)$ it continues for one
forcing period, 100 steps, and the hold is the fraction of those steps
spent within 0.06 of $\mathbf{p}^*_1(t)$. At $\epsilon = 0$ both trackers
succeed, keep the straddle, and hold in every trial, as in
Section~\ref{sec:disc_clean}.

\begin{table}[H]
\centering
\caption{Time-periodic double gyre, noise-free, $10^3$ trials per start
pooled over the seven starts. Rates in percent. Error is the median over
trials of the mean $|x_c - x_s(t)|$ after acquisition, and hold is the
median over successful trials. Error is not reported at
$\epsilon = 0.25$, where only 35\% ($D$) and 15\% ($s_1$) of trials
acquire the separatrix.}
\label{tab:unsteady}
\begin{tabular}{llrrrrr}
\toprule
$\epsilon$ & Tracker & Success & Retention & Reversal & Error & Hold \\
\midrule
0    & $D$   & 100  & 100  & 0    & 0.004 & 1.00 \\
0    & $s_1$ & 100  & 100  & 0    & 0.010 & 1.00 \\
0.1  & $D$   & 97.3 & 97.3 & 2.7  & 0.017 & 1.00 \\
0.1  & $s_1$ & 96.2 & 59.1 & 3.8  & 0.025 & 1.00 \\
0.25 & $D$   & 30.0 & 8.6  & 16.1 & n/a   & 0.39 \\
0.25 & $s_1$ & 59.4 & 15.5 & 28.0 & n/a   & 0.27 \\
\bottomrule
\end{tabular}
\end{table}

At $\epsilon = 0.1$ both trackers reach $\mathbf{p}^*_1(t)$ in more than
96\% of trials and hold it for the full period
(Table~\ref{tab:unsteady}), since the saddle moves along the wall at no
more than 0.063, below the $k\,c_{\max} = 0.12$ that each component of
either hold law can command. The $D$ tracker keeps the straddle in every
trial that succeeds, and its median tracking error of 0.017 is close to
the 0.015 to 0.019 offset of the true $D$ trench from the separatrix at
peak forcing. The $s_1$ tracker reaches the saddle as often but keeps the
straddle in only 59\% of trials, and in 9 to 14\% from S1, S2, and S3. Of
the successful trials that lose it, 97\% lose it within 0.1 of $y = 0$.
At $y = 0$ the strain on the separatrix falls to at most
$\pi A\epsilon = 0.031$, against $\pi^2 A = 0.99$ at the saddles, so $s_1$
is nearly flat across the trench there, and the across-trench command of
(\ref{eq:s1_law}), which scales with the slope, cannot keep pace with the
moving trench. The $D$ tracker's across-trench command is a Newton step,
a distance, which the flattening does not weaken.

Reversals at $\epsilon = 0.1$ occur only from S1 and S2, and for the
$s_1$ tracker also from S4, in 7 to 11\% of their trials. None of them
rides down first. Each starts at a phase where the displaced separatrix
is nearly as far from the start as the top-wall trench, with a median
ratio of the two distances of 0.98, against 0.47 in trials that do not
reverse. The cluster acquires the wrong trench rather than turning on the
right one.

At $\epsilon = 0.25$ the separatrix moves at up to 0.157, faster than the
0.12 cap on either tracker's across-trench channel, and neither tracker
keeps up. The $D$ tracker succeeds in 30\% of trials and the $s_1$
tracker in 59\%, and neither holds the saddle, which outruns both hold
laws along the wall. In 91\% of the $D$ tracker's reversals and 86\% of
the $s_1$ tracker's, the top-wall trench was nearer to the start than the
displaced separatrix. The $D$ tracker's other failures, 54\% of its
trials, miss the moving saddle and leave the basin at a bottom corner.
With $c_{\max}$ doubled to 0.08, which raises the cap to 0.24, the $D$
tracker succeeds in 62\% of trials and holds for a median 0.89 of the
period, and the $s_1$ tracker in 62\% with a full hold. Reversals remain at
18\% and 26\%, since they come from the start geometry rather than the
speed limit.

\subsection*{IV-F \quad Bent Double Gyre}

The time-periodic gyre keeps a straight separatrix. A bent gyre curves it
while staying steady and incompressible, with stream function
$\psi = A\sin(\pi g)\sin(\pi y_f)$, $g = x_f + \delta\sin(\pi y_f)$,
$u = -\partial\psi/\partial y$, and $v = \partial\psi/\partial x$. At
$\delta = 0$ it is (\ref{eq:dg_u}). For $\delta > 0$ the separatrix is the
curve $x = -\delta\cos(\pi y)$, bowed by $\delta$ at $y = 0$, between the
same saddles, and the flow on it still runs downward. On it the vorticity
reaches 0.60 at $\delta = 0.1$ and 1.19 at $\delta = 0.2$, comparable to
the strain $\pi^2 A = 0.99$ at the saddles. Neither surrogate's trench
follows the curve. Both bow further than the separatrix and meet it near
the saddles, the trench of $D$ by up to 0.041 and 0.078 and the trench of
$s_1$ by up to 0.085 and 0.155, so at mid-height the $s_1$ trench lies
more than $\rho$ from the separatrix.

\begin{table}[H]
\centering
\caption{Bent double gyre, noise-free, 500 trials per start. Success is
split between the starts at or above the origin (S1 to S4) and below it
(S5 to S7). Retention and reversal are pooled over all seven starts.
Rates in percent.}
\label{tab:bent}
\begin{tabular}{llrrrr}
\toprule
$\delta$ & Tracker & S1--S4 & S5--S7 & Retention & Reversal \\
\midrule
0   & $D$   & 100  & 100  & 100  & 0   \\
0   & $s_1$ & 100  & 100  & 100  & 0   \\
0.1 & $D$   & 100  & 33.3 & 52.5 & 0   \\
0.1 & $s_1$ & 99.6 & 100  & 23.5 & 0.3 \\
0.2 & $D$   & 100  & 11.5 & 9.7  & 0   \\
0.2 & $s_1$ & 84.6 & 11.9 & 0.5  & 9.6 \\
\bottomrule
\end{tabular}
\end{table}

From the starts at or above the origin, the $D$ tracker reaches
$\mathbf{p}^*_1$ in every trial at both bends, and the $s_1$ tracker in
99.6\% at $\delta = 0.1$ and 84.6\% at $\delta = 0.2$
(Table~\ref{tab:bent}). Both hold it in every successful trial at
$\delta = 0.1$. Straddle retention falls for both, and furthest for the
$s_1$ tracker, because each follows its own trench rather than the
separatrix. Its median offset from the separatrix after acquisition is
0.048 for the $D$ tracker and 0.075 for the $s_1$ tracker at
$\delta = 0.1$. From S6 and S7 the $D$ tracker fails in every trial at
both bends, leaving the basin at its bottom-right corner, while the $s_1$
tracker succeeds from all three lower starts at $\delta = 0.1$.

\subsection*{IV-G \quad Bickley Jet}

Both gyres have separatrices on which strain dominates over most of their
length. A Bickley jet seen in the frame of its wave,
$\psi = -UL\tanh(y/L) + cy + UL\epsilon\,\mathrm{sech}^2(y/L)\cos k(x + 0.5)$,
has a row of cat's-eye eddies on each flank, each bounded by curved
separatrices that run from saddle to saddle inside a sheared jet. With
$U = 0.3$, $L = 0.2$, $c = 0.3U$, $\epsilon = 0.3$, and $k = 2\pi$, the
lower separatrix of one northern eye runs east from the saddle at
$(-0.5, 0.29)$ to the saddle at $(0.5, 0.29)$, dipping 0.26 below them,
with length 1.14 and an eye 0.28 high. On it the vorticity ranges from
0.82 to 1.96, against a strain rate of 0.33 at the saddles. The trench of
$s_1$ lies within 0.13 of the separatrix, and $D$ has no trench within
0.15 of it at 40\% of positions along it.

Neither tracker follows this separatrix. Over $10^3$ trials from each of
seven starts, three on the separatrix and four off it by $1.3\rho$ or
$2\rho$, the $D$ tracker reaches the downstream saddle in 4.1\% of trials,
all from the two starts nearest it, and the $s_1$ tracker in none, and
neither keeps the straddle in any trial. The $D$ tracker leaves the
domain in 56\% of trials and reaches the upstream saddle in 9\%. The $s_1$
tracker runs the full 500 steps without reaching either saddle in 99.7\%
of trials, ending on its own trench, which runs nearly flat along the jet
flank below the saddles. Both trackers steer on the trenches of their
surrogate fields, and on a separatrix where rotation outweighs strain
those trenches do not mark it.

\section{Edits to existing sentences of Draft 13}

Each pair below quotes the current text exactly as it stands in
Draft~13.tex and gives the proposed replacement. None is applied.

\paragraph{1. Section~\ref{sec:surrogates}, Determinant Field.}
\emph{Current:} ``In the flows studied here, separatrices run along the
trenches of $D$ inside the strain-dominated regions.''\\
\emph{Proposed:} ``A separatrix runs along a trench of $D$ where the
vorticity on it is small against the strain, as on the steady double
gyre, where the vorticity on the separatrix vanishes.''\\
\emph{Ripple:} Section~\ref{sec:dg_example} already states the vanishing
vorticity, so nothing there changes.

\paragraph{2. Abstract.}
\emph{Current:} ``Each tracker locates the nearest separating boundary
and navigates to and along it from instantaneous measurements, with no
integration horizon or pre-straddling requirements.''\\
\emph{Proposed:} ``Each tracker locates the nearest trench of its derived
field and navigates to and along it from instantaneous measurements, with
no integration horizon or pre-straddling requirements. These trenches
mark the separating boundary where strain outweighs rotation on it.''\\
\emph{Ripple:} the preceding sentence says the primitives steer ``along
the transport boundaries,'' which reads as a general claim; it could
become ``toward the transport boundaries.''

\paragraph{3. Section~\ref{sec:disc_limits}, first $s_1$ sentence.}
\emph{Current:} ``The $s_1$ tracker's tangent selection weakens near the
saddles, and its passage through the origin rests on the benchmark's
vanishing shear strain.''\\
\emph{Proposed:} ``The $s_1$ tracker's tangent selection weakens near the
saddles. On the time-periodic gyre, where the shear strain no longer
vanishes, it still passes the crest, but at $\epsilon = 0.1$ it reaches the
saddle without keeping the straddle in 37\% of trials, and 97\% of those
lose it within 0.1 of the crest, where $s_1$ flattens across the
trench.''\\
\emph{Ripple:} the next sentence begins ``Without it the strain
eigenvector can tilt off the trench,'' and its ``it'' would need a new
antecedent.

\paragraph{4. Section~\ref{sec:disc_limits}, new paragraph after the
evaluation paragraph.}
\emph{Proposed:} ``Both trackers follow trenches of their surrogate
fields, and these mark the separatrix only where strain outweighs
rotation on it. On the time-periodic gyre, where the median
rotation-to-strain ratio $|\omega|/2r$ along the separatrix is 0.018 at
$\epsilon = 0.1$, the $D$ tracker keeps the straddle in 97\% of trials. On
the bent gyre the median rises to 0.25 and 0.46 and retention falls to
53\% and 10\%, and on the Bickley jet, at 0.79, neither tracker keeps it.
The cluster measures the same ratio from the first-order coefficients of
its fit, and its median along the Santa Barbara paths was 0.15 to 0.20. A
cluster could therefore flag online when the trench it rides is unlikely
to mark a separatrix.''\\
\emph{Ripple:} the 37\% and 97\% figures depend on IV-E being accepted,
and the bent and Bickley figures on IV-F and IV-G.

%======================================================================
\appendix
\part*{Appendices}
\addcontentsline{toc}{part}{Appendices}
%======================================================================

\section{Provenance}

All scripts are under
\begin{center}\small\path{trunk/Python_Simulations/Vector_Fields/VF_Robot/experiments/unsteady_dg/}\end{center}
with outputs in its \texttt{outputs/} folder, run from \texttt{VF\_Robot}
with \texttt{venv/bin/python3}. Nothing is committed.

\begin{center}
\small
\begin{tabular}{p{0.36\textwidth}p{0.58\textwidth}}
\toprule
Command & Produces \\
\midrule
\texttt{mc\_unsteady\_double\_gyre.py --trials 1000 --eps 0 0.1 0.25 --tag n1000} & \texttt{trials\_n1000.csv}, 42\,000 trials, 500\,s \\
\texttt{mc\_unsteady\_double\_gyre.py --trials 1000 --eps 0.25 --vmax 0.08 --tag eps025\_vmax008} & the $2c_{\max}$ diagnostic, 14\,000 trials \\
\texttt{mc\_bent\_double\_gyre.py --trials 500 --tag n500} & \texttt{bent\_n500.csv}, 21\,000 trials, 258\,s \\
\texttt{mc\_bent\_double\_gyre.py --diag} & bent-gyre vorticity and trench offsets \\
\texttt{mc\_bickley\_jet.py --trials 1000 --tag n1000} & \texttt{bickley\_n1000.csv}, 14\,000 trials, 372\,s \\
\texttt{mc\_bickley\_jet.py --geom} & jet geometry, separatrix checks \\
\texttt{trench\_offset.py} & time-periodic amplitudes, speeds, offsets \\
\texttt{scope\_ratio.py} & Table~\ref{tab:ratio}, true columns \\
\texttt{path\_ratio.py --trials 50 --tag n50} & Table~\ref{tab:ratio}, fitted columns \\
\texttt{analyze.py --tag n1000} & per-start, phase, reversal, and exit summaries \\
\texttt{fig\_notes.py} & Figs.~\ref{fig:structures}, \ref{fig:crest}, \ref{fig:rides}, and \ref{fig:scope}, and the fine-sampled trench offsets of Table~\ref{tab:fields} \\
\bottomrule
\end{tabular}
\end{center}

The fitted-ratio reruns wrap both primitives in each worker process to log
the fit before each step. The primitives' code is unchanged, and the
reruns reproduce the full runs' success and retention to within sampling
error (for example 98.0\% against 97.3\% for the \Dtr{} at
$\epsilon = 0.1$). The trench offsets in Table~\ref{tab:fields} come from
fine sampling (197 heights or positions) and supersede coarser values
quoted earlier in the session (0.093 for the $s_1$ trench on the jet, and
7 of 17 positions without a $D$ trench).

\section{Edits applied to Draft 13 earlier in this session}

These are in \texttt{Draft\_13.tex} now. The paper compiles at 11 pages
with the same warnings as before the edits.
\begin{enumerate}
\item Section~\ref{sec:dg_bench}: the task is stated as riding the
separatrix downstream to $\mathbf{p}^*_1$, a run ending at
$\mathbf{p}^*_2$ fails, and the straddle is defined as robots on both
sides of $x = 0$.
\item ``Far saddle'' became ``downstream saddle'' in all seven places.
\item Section~\ref{sec:disc_ocean}: the $s_1$ typing now reports the
eigenvector followed and notes that the forward FTLE reference cannot
confirm the attracting segment.
\item Section~\ref{sec:arch} and the caption of
Fig.~\ref{fig:control_architecture} name the reused architecture and the
two blocks new in this paper; the formation cites \cite{5}.
\item To hold 11 pages, Fig.~\ref{fig:pentagon} is drawn at 0.8 and
Fig.~\ref{fig:dg_fields} at 0.95 of the column width.
\end{enumerate}
Still open from that session: Fig.~\ref{fig:flip_resolution} marks the
50\% crossings of the success curves while the text and abstract compare
retention, and from S6 and S7 the two point opposite ways; the scoring
code also fails a run whose formation collapses (pair-distance RMS error
above 0.30), which the paper does not mention; and \texttt{Ready to
Submit/} is out of date.

% Keys are display-only: \cite{N} resolves to Draft_13's [N] through xr,
% which imports its \bibcite entries.
\nocite{list5,list8,list12,list17,list19,list22}
\begin{thebibliography}{99}
\bibitem[5]{list5} L. Bri\~{n}\'{o}n-Arranz, A. Renzaglia, and L. Schenato,
``Multirobot symmetric formations for gradient and Hessian estimation with
application to source seeking,'' \textit{IEEE Trans. Robot.}, vol. 35,
no. 3, pp. 782--789, 2019.
\bibitem[8]{list8} A. Okubo, ``Horizontal dispersion of floatable particles in
the vicinity of velocity singularities such as convergences,''
\textit{Deep-Sea Res.}, vol. 17, no. 3, pp. 445--454, 1970.
\bibitem[12]{list12} M. Serra et al., ``Search and rescue at sea aided by
hidden flow structures,'' \textit{Nat. Commun.}, vol. 11, no. 1, p. 2525,
2020.
\bibitem[17]{list17} S. C. Shadden, F. Lekien, and J. E. Marsden,
``Definition and properties of Lagrangian coherent structures from
finite-time Lyapunov exponents in two-dimensional aperiodic flows,''
\textit{Physica D}, vol. 212, no. 3-4, pp. 271--304, 2005.
\bibitem[19]{list19} M. Serra and G. Haller, ``Objective Eulerian coherent
structures,'' \textit{Chaos}, vol. 26, no. 5, p. 053110, 2016.
\bibitem[22]{list22} J. Weiss, ``The dynamics of enstrophy transfer in
two-dimensional hydrodynamics,'' \textit{Physica D}, vol. 48, no. 2-3,
pp. 273--294, 1991.
\end{thebibliography}

\end{document}
